\subsection{Conceptual Foundations}

\begin{enumerate}[leftmargin=*, label=\textbf{Q\arabic*.}]

\item State the \textbf{Implicit Function Theorem} (IFT) precisely. What are the hypotheses on $F : \mathbb{R}^{n+m} \to \mathbb{R}^m$, on the point $(a, b) \in \mathbb{R}^n \times \mathbb{R}^m$, and what does the conclusion guarantee?

\item In the IFT, why is it necessary to require that the \textbf{partial Jacobian} $D_y F(a,b)$ is invertible (i.e.\ $\det D_y F(a,b) \neq 0$)? What can go wrong if this condition fails? Provide a concrete example.

\item The IFT guarantees a function $g : U \to V$ such that $F(x, g(x)) = 0$ for all $x \in U$. What are $U$ and $V$? In what sense is $g$ \textbf{unique}?

\item Explain the relationship between the \textbf{Implicit Function Theorem} and the \textbf{Inverse Function Theorem}. Can each be derived from the other?

\item The IFT is a \textbf{local} result. What does ``local'' mean precisely here, and why can the conclusion generally not be extended to a global statement?

\item What regularity does the IFT guarantee for $g$? If $F \in C^k$, what can you conclude about $g$?

\end{enumerate}

\subsection{Computing the Implicit Derivative}

\begin{enumerate}[resume, leftmargin=*, label=\textbf{Q\arabic*.}]

\item Derive the formula for the derivative $Dg(x)$ in terms of the partial derivatives of $F$. Starting from the identity $F(x, g(x)) = 0$, apply the chain rule and solve for $Dg(x)$.

\item In the scalar case $F(x, y) = 0$ with $F : \mathbb{R}^2 \to \mathbb{R}$, write down the explicit formula
\[
\frac{dy}{dx} = -\frac{\partial F / \partial x}{\partial F / \partial y}.
\]
Derive this from the chain rule and identify when the formula is valid.

\item Let $F(x, y) = x^2 + y^2 - 1$. Near which points $(a,b)$ on the unit circle does the IFT guarantee a local function $y = g(x)$? Near which points does it \emph{fail}, and why?

\item Consider $F(x,y,z) = x^2 + y^2 + z^2 - 1 = 0$ and suppose we wish to express $z$ as a function of $(x,y)$. State the condition under which this is possible and compute $\frac{\partial z}{\partial x}$ and $\frac{\partial z}{\partial y}$ implicitly.

\item Let $F : \mathbb{R}^2 \times \mathbb{R}^2 \to \mathbb{R}^2$ be defined by
\[
F(x_1, x_2, y_1, y_2) = \begin{pmatrix} x_1^2 - y_1 + y_2 \\ x_2 + y_1 y_2 - 1 \end{pmatrix}.
\]
At the point $(0,1,0,1)$, verify the hypothesis of the IFT and write down the system you would solve to find $Dg(0,1)$.

\end{enumerate}

\subsection{Geometric and Analytic Interpretation}

\begin{enumerate}[resume, leftmargin=*, label=\textbf{Q\arabic*.}]

\item Interpret the IFT geometrically: what does it say about the \textbf{level set} $\{(x,y) : F(x,y) = 0\}$ near a regular point?

\item Define a \textbf{regular point} and a \textbf{critical point} of $F : \mathbb{R}^{n+m} \to \mathbb{R}^m$. How do these concepts relate to the hypothesis of the IFT?

\item The IFT implies that the zero set of $F$ is locally a \textbf{smooth manifold} near a regular point. What is the dimension of this manifold, and why?

\item Explain how the IFT is used to verify that a subset $M \subseteq \mathbb{R}^n$ defined by $k$ independent equations is a \textbf{smooth submanifold} of dimension $n - k$.

\end{enumerate}

\subsection{Applications and Extensions}

\begin{enumerate}[resume, leftmargin=*, label=\textbf{Q\arabic*.}]

\item State the \textbf{rank theorem} (constant rank theorem). How does it generalise the IFT?

\item In constrained optimisation, the method of \textbf{Lagrange multipliers} relies on the IFT. Explain what role the IFT plays in justifying why the constraint surface can be locally parameterised, thus reducing the problem.

\item Let $F(x, \lambda) = 0$ define an \textbf{implicitly parameterised curve} in $(x,\lambda)$-space. What does the IFT say about the local structure of this curve? Where might the curve \textbf{bifurcate}, and how is this detected?

\item Suppose $F(x,y) = 0$ and $G(x,y) = 0$ are two equations in two unknowns. State conditions under which the system has a locally unique solution $(x,y) = (a,b)$ near a given solution, and relate this to the IFT.

\item In the context of \textbf{ordinary differential equations}, explain how the IFT is used to prove a local \textbf{existence and uniqueness theorem} for solutions of $\dot{x} = f(x,t)$.

\item Consider the map $\Phi : \mathbb{R}^n \to \mathbb{R}^n$ and a point where $\Phi(a) = b$. State the \textbf{Inverse Function Theorem} and show how it follows from the IFT applied to $F(x,y) = \Phi(y) - x$.

\end{enumerate}

\subsection{Proof and Technique}

\begin{enumerate}[resume, leftmargin=*, label=\textbf{Q\arabic*.}]

\item Sketch the main idea of the \textbf{proof of the IFT}. Which fixed-point theorem (or contraction principle) is typically invoked, and how?

\item In the proof of the IFT via the \textbf{Banach fixed-point theorem}, one defines an auxiliary map $T_x(y) = y - [D_y F(a,b)]^{-1} F(x,y)$. Explain why $T_x$ is a contraction on a suitable ball, and what the fixed point represents.

\item Why does the proof of the IFT require $F$ to be at least $C^1$? Give an example of a function that satisfies $F(0,0) = 0$ and $\partial F/\partial y \neq 0$ at the origin but for which the conclusion of the IFT fails due to lack of continuity of derivatives.

\item After establishing existence and uniqueness of $g$ from the IFT, how do you then prove that $g$ is \textbf{differentiable} (not just continuous)? Outline the key steps.

\end{enumerate}

\subsection{Quick-Fire Recall}

\begin{enumerate}[resume, leftmargin=*, label=\textbf{Q\arabic*.}]

\item True or False: if $F(a,b) = 0$ and $D_y F(a,b)$ is invertible, then $g$ is defined on all of $\mathbb{R}^n$.

\item True or False: the IFT guarantees that $g$ is real-analytic whenever $F$ is real-analytic.

\item What is the dimension of the solution set of $m$ independent equations $F(x) = 0$ in $\mathbb{R}^n$ (with $m < n$) near a regular point?

\item If $F : \mathbb{R}^3 \to \mathbb{R}$ and $\nabla F(a) \neq 0$, what geometric object does the IFT say the level surface $F = 0$ looks like near $a$?

\item Write the formula for $\frac{\partial(y_1, y_2)}{\partial(x_1, x_2)}$ in terms of partial derivatives of $F = (F_1, F_2)$ using the IFT.

\end{enumerate}
